\documentclass[10pt]{article}
\usepackage[T1]{fontenc}
\usepackage[margin=0.9in]{geometry}
\usepackage{amsmath,amssymb,mathtools}
\usepackage{booktabs,array}
\usepackage{xcolor}
\usepackage[hidelinks]{hyperref}
\usepackage{microtype}
\usepackage{fancyhdr}
\definecolor{navy}{RGB}{24,55,95}
\definecolor{green}{RGB}{30,110,70}
\pagestyle{fancy}
\setlength{\headheight}{14pt}
\fancyhf{}
\lhead{Finite linear groups}
\rhead{$\sigma$-MGF proof and exact verification}
\cfoot{\thepage}
\newcommand{\Sym}{\operatorname{Sym}}
\newcommand{\Fix}{\operatorname{Fix}}
\newcommand{\F}{\mathbb F}
\newcommand{\Sg}{\mathcal S}
\title{\color{navy}\textbf{Finite Linear-Group Permutation Actions}\\[2mm]
\large Exact $\sigma$-MGFs, stable coefficients, and matrix verification}
\author{Proof and computational verification companion}
\date{17 July 2026}

\begin{document}
\maketitle

\begin{abstract}
For a finite group acting on a set $S$, we prove the determinant/cycle formula,
its fixed-point reconstruction, and its interpretation as an orbit count on
tuples of multisets.  We then specialize to nonzero-vector, projective-point,
and complete-flag actions of finite linear groups.  Exact finite-field matrix
enumerations verify all coefficients through degree three for
$PSL_3(2)$, the degree-two vector claims for $GL_2(3)$, and the Bruhat
obstruction in ranks two and three.  The computation also exposes a sharp
boundary phenomenon for $SL_2(3)$.
\end{abstract}

\section{The universal finite permutation formula}
Let $G$ act on a finite set $S$, let $V=\mathbb C[S]$, and write $P_g$ for
the associated permutation matrix.  For a partition
$\tau=(\tau_1,\ldots,\tau_s)$ set
\[
 W_\tau=\bigotimes_{b=1}^s\Sym^{\tau_b}(V).
\]
The symmetric-power identity
\[
 \det(I-tP_g)^{-1}=\sum_{r\geq0}
 \operatorname{tr}\!\left(P_g\mid\Sym^rV\right)t^r
\]
implies, after averaging over $G$,
\begin{equation}
 [m_\tau]\Sg_{G,S}
 =\frac1{|G|}\sum_{g\in G}\operatorname{tr}(g\mid W_\tau)
 =\dim W_\tau^G.
\end{equation}
The last equality follows because $|G|^{-1}\sum_g g$ is the Reynolds
projector.

If $g$ has $c_d(g;S)$ cycles of length $d$, a single $d$-cycle has all
$d$th roots of unity as eigenvalues and hence determinant $1-t^d$.
Therefore
\begin{equation}
 \boxed{\Sg_{G,S}(t_\bullet)=\frac1{|G|}\sum_{g\in G}
 \prod_{i\geq1}\prod_{d\geq1}(1-t_i^d)^{-c_d(g;S)}.}
\end{equation}
This is an exact finite formula, not an asymptotic statement.

\subsection*{Fixed points determine the cycles}
A point on a $d$-cycle is fixed by $g^r$ exactly when $d\mid r$.  Thus
\[
 \Fix_S(g^r)=\sum_{d\mid r}d\,c_d(g;S).
\]
M\"obius inversion gives
\begin{equation}
 \boxed{c_d(g;S)=\frac1d\sum_{r\mid d}
 \mu(d/r)\Fix_S(g^r).}
\end{equation}
For $S=G/P$, $\Fix_S(g^r)$ is the value at $g^r$ of the permutation
character $\operatorname{Ind}_P^G\mathbf1$.  Grouping (2) by conjugacy class
then yields the stated centralizer-weighted formula.

\subsection*{Direct orbit interpretation}
A monomial basis vector of $\Sym^r\mathbb C[S]$ is a multiset of $r$ points
of $S$.  Consequently a basis of $W_\tau$ is a tuple of such multisets, of
respective sizes $\tau_1,\ldots,\tau_s$.  Since $G$ permutes this basis, the
sum of the basis vectors in each orbit is one invariant basis vector.  Hence
\begin{equation}
 \boxed{[m_\tau]\Sg_{G,S}=\left|\{G\text{-orbits on tuples of multisets of
 sizes }\tau_1,\ldots,\tau_s\}\right|.}
\end{equation}
Equation (4) supplies the independent direct route used in the computation.

\section{Linear and projective specializations}
For $S=\F_q^n\setminus\{0\}$ and $g\in GL_n(q)$,
\begin{equation}
 \Fix_S(g^r)=q^{\dim\ker(g^r-I)}-1.
\end{equation}
For projective points, a fixed line is an eigenline whose eigenvalue lies in
$\F_q^\times$.  The eigenspaces for distinct eigenvalues contribute disjoint
sets of lines, so
\begin{equation}
 \Fix_{\mathbb P^{n-1}(\F_q)}(g^r)=
 \sum_{\lambda\in\F_q^\times}
 \frac{q^{\dim\ker(g^r-\lambda I)}-1}{q-1}.
\end{equation}
Substitution of (5) or (6) into (3) and then (2) proves the proposed exact
formulas.

Over $\F_2$, every projective line has one nonzero vector and
$GL_n(2)=SL_n(2)=PGL_n(2)=PSL_n(2)$.  Thus (5) directly gives the displayed
$PSL_n(2)$ formula.

\subsection*{Stable range}
A configuration of total degree $D=|\tau|$ involves at most $D$ projective
points and its span has dimension at most $D$.  When $n\geq D$, two such
configurations are $GL_n(q)$-equivalent exactly when the induced colored
configurations on their spans are linearly isomorphic: any isomorphism of the
spans extends after choosing complements.  This proves coefficientwise
stability for $GL/PGL$ projective actions.  The same argument for nonzero
vectors is immediate.  For $SL/PSL$, the determinant of an extension can be
corrected on an unused one-dimensional summand, which requires the strict
range $n>D$.

For projective points the stable expansion through degree three is
\begin{equation}
 1+m_{(1)}+2m_{(2)}+2m_{(1,1)}
 +4m_{(3)}+5m_{(2,1)}+6m_{(1,1,1)}+O(\deg4).
\end{equation}
For three distinct points, collinear and noncollinear triples give two
different orbits; this is the extra degree-three term relative to the
symmetric-group/Poisson pattern.

For nonzero vectors, ordered pairs consist of one independent orbit and one
dependent orbit $y=ax$ for each $a\in\F_q^\times$, giving
\[
 [m_{(1,1)}]\Sg=q.
\]
For an unordered pair, the dependent ratios other than $1$ are identified by
$a\leftrightarrow a^{-1}$.  Together with the repeated and independent types,
this gives
\[
 [m_{(2)}]\Sg=2+\left\lfloor\frac{q-1}{2}\right\rfloor.
\]

\section{Complete flags and failure of rank stability}
For $S=G/B$, orbitals are the double cosets $B\backslash G/B$.  Bruhat
decomposition identifies these with the Weyl group, whence
\begin{equation}
 [m_{(1,1)}]\Sg_{G,G/B}=|W|.
\end{equation}
In type $A_{n-1}$ this is $n!$, so even the degree-two coefficient cannot
stabilize as rank grows.

\section{Exact verification method}
The accompanying marimo notebook performs the following independent steps.
\begin{enumerate}
 \item Enumerate every finite-field matrix and retain the required determinant
 or invertibility condition.
 \item Induce its permutation on nonzero vectors, normalized projective
 representatives, or incident point--hyperplane flags.
 \item Count orbits directly on tuples of multisets, without using character
 or cycle formulas.
 \item Extract the same coefficient from the actual cycle counts in (2).
 \item Reconstruct every cycle count from fixed points of powers using (3).
 \item At $(t_1,t_2)=(0.07,0.11)$, compare the direct matrix-determinant
 average with the cycle-product average.
\end{enumerate}
The last two numerical averages both equal $1.247556190416$ for the
$PSL_3(2)$ projective action, with difference below displayed precision.

\begin{center}
\small
\begin{tabular}{@{}lcrrr@{}}
\toprule
Action & $\tau$ & direct orbits & cycle formula & claimed\\
\midrule
$PSL_3(2)$ on $\mathbb P^2(\F_2)$ & $(1)$ & 1&1&1\\
 & $(2)$ &2&2&2\\
 & $(1,1)$ &2&2&2\\
 & $(3)$ &4&4&4\\
 & $(2,1)$ &5&5&5\\
 & $(1,1,1)$ &6&6&6\\
$PSL_2(2)$ on $\mathbb P^1(\F_2)$ & $(3)$ &3&3&3\\
$GL_2(3)$ on nonzero vectors & $(2)$ &3&3&3\\
 & $(1,1)$ &3&3&3\\
$SL_2(3)$ on nonzero vectors & $(1,1)$ &4&4&4\\
$PGL_2(3)$ on projective points & $(1,1)$ &2&2&2\\
$PSL_2(3)$ on projective points & $(1,1)$ &2&2&2\\
$GL_2(2)$ on complete flags & $(1,1)$ &2&2&2\\
$GL_3(2)$ on complete flags & $(1,1)$ &6&6&6\\
\bottomrule
\end{tabular}
\end{center}

The constructed orders are $|GL_3(2)|=168$, $|GL_2(3)|=48$, and
$|SL_2(3)|=24$; their projective images in degree two have orders 24 and 12.
All displayed tests pass exactly.

\paragraph{Sharpness found by verification.}
The $SL_2(3)$ value $4$ at $\tau=(1,1)$ differs from the $GL_2(3)$ value
$3$.  Here $n=|\tau|=2$, so there is no unused complement on which to correct
the determinant.  Independent ordered pairs split into two $SL_2(3)$ orbits
according to their determinant.  This confirms that the proposed $SL/GL$
agreement must retain the strict hypothesis $n>|\tau|$.

\section{Scope of the conclusion}
The matrix, orbit, fixed-point, and determinant routes agree for every
representative case above.  These computations verify the exact mechanism and
the low-degree stable claims; the proofs of (1)--(8), rather than finite
sampling, establish the formulas for all finite parameters in their stated
ranges.

\end{document}
